%% 5-10 細い長方形に分けてたそう（第11回：ページ制限なし。X16・X18・X19 に対応して書き足した）
\bkchapter{5}{10}{細い長方形に分けてたそう}{\bklevelrange{2}{4}}{細い長方形に分けてたそう}
\bkmichishirube{4-5,5-3,5-9}{5-11,5-12}

\begin{goal}
\begin{itemize}[label={\color{bkTheme}●}, leftmargin=1.6\zw, itemsep=0pt, topsep=1pt]
\item \bklevel{2}\enspace 曲線の下の面積を、細い長方形の和で近似できる。
\item \bklevel{3}\enspace $\Sigma$ の記号を読み書きできる。
\item \bklevel{4}\enspace 定積分 $\displaystyle\int_a^b f(x)\,\dd x$ の意味を説明し、$x^2$ の定積分を求められる。
\end{itemize}
\end{goal}

\begin{hatena}
\noindent\begin{minipage}[c]{0.56\linewidth}
落ちるボールの速さは $v=10t$ で、グラフは直線でした。だから、進んだ距離は三角形の面積で求められました（5-9）。

では、速さが $v=t^2$ のように、グラフが曲がっているときはどうでしょう。右の色をぬった部分（0秒から1秒までに進んだ距離）は、三角形でも台形でもありません。どうやって面積を求めればよいでしょう。
\end{minipage}\hfill
\begin{minipage}[c]{0.4\linewidth}
\begin{tikzpicture}
\begin{axis}[bkplot, width=54mm, height=42mm, xmin=0, xmax=1.25, ymin=0, ymax=1.2, xtick={1}, ytick={1}, xlabel={$t$}, ylabel={$v$}]
  \addplot[fill=bkYellow!40, draw=none, domain=0:1] {x^2} \closedcycle;
  \addplot[bkcurve, domain=0:1.1] {x^2};
  \node[font=\sffamily\bfseries\small] at (axis cs:0.72,0.2) {?};
\end{axis}
\end{tikzpicture}
\end{minipage}
\end{hatena}

5-9では、三角形の面積を「細い長方形の階段」で考えました。曲がったグラフでも、同じ方法が使えます。細い長方形に分けて、たし合わせるのです。この「細かく分けてたす」計算が、\textbf{積分}です。

この章は、微分の章（5-1〜5-3）と同じように、1段ずつ進みます。はじめに方眼のマスを数え、次に長方形を4本だけならべて手で計算し、それから本数をふやします。最後に、本数をかぎりなくふやした先に名前をつけます。

%% ---------------- やってみよう・思い出そう ----------------
\begin{yattemiyou}[方眼のマスを数えて、面積を見積もろう]
はてな？の $v=t^2$ を、数学でよく使う文字で $y=x^2$ と書きます（横軸が $x$、たて軸が $y$。グラフの形は同じです）。右の方眼は、たても横も1目もりが0.1です。$x=0$ から1までの、グラフの下の部分のマスを数えます。

\noindent\begin{minipage}[c]{0.44\linewidth}\raggedright
\begin{enumerate}[label={\bkcircled[bkGreen]{\arabic*}}, leftmargin=2.4\zw, itemsep=2pt, topsep=3pt]
\item 全部がグラフの下に入っているマス（青）は、何こかな？
\item グラフが通っていて、一部だけ入っているマス（黄）は、何こかな？
\item 黄のマスは、2つで1マスと数えよう（5-9と同じ）。全部で何マスかな？
\item 1マスの面積は $0.1\times0.1=0.01$。グラフの下の面積は、およそいくつかな？
\end{enumerate}
\end{minipage}\hfill
\begin{minipage}[c]{0.53\linewidth}\centering
\begin{tikzpicture}[x=48mm, y=48mm]
  % 第11回：列ごとに、全部入るマスの数（青）と、一部だけ入るマスの数（黄）。Python で数えた
  \foreach \i/\f/\p in {1/0/1,2/0/1,3/0/1,4/0/2,5/1/2,6/2/2,7/3/2,8/4/3,9/6/3,10/8/2} {
    {\ifnum\f>0 \fill[bkBlue!35] ({(\i-1)/10},0) rectangle ({\i/10},{\f/10});\fi
     \fill[bkYellow!60] ({(\i-1)/10},{\f/10}) rectangle ({\i/10},{(\f+\p)/10});}}
  \draw[bkGray!70, line width=0.3pt] (0,0) grid[step=0.1] (1,1);
  \draw[bkBlue!80!black, line width=1.3pt] plot[domain=0:1, samples=60] (\x,{\x*\x});
  \draw[bkInk, line width=0.7pt, -{Stealth[length=4pt]}] (0,0) -- (1.09,0) node[right, font=\sffamily\small, inner sep=1pt] {$x$};
  \draw[bkInk, line width=0.7pt, -{Stealth[length=4pt]}] (0,0) -- (0,1.09) node[above, font=\sffamily\small, inner sep=1pt] {$y$};
  \foreach \v in {0.5,1} {
    \node[font=\sffamily\footnotesize, anchor=north, inner sep=1pt] at (\v,-0.015) {\v};
    \node[font=\sffamily\footnotesize, anchor=east, inner sep=1pt] at (-0.015,\v) {\v};}
  \node[font=\sffamily\footnotesize, anchor=north east, inner sep=1pt] at (-0.01,-0.01) {0};
  \node[font=\sffamily\footnotesize, text=bkBlue!80!black, anchor=west, inner sep=1pt] at (1.02,0.93) {$y=x^2$};
\end{tikzpicture}
\end{minipage}
\end{yattemiyou}

\noindent 青のマスは24こ、黄のマスは19こです。黄は2つで1マスと数えるので、$19\div2=9.5$ マス。全部で $24+9.5=33.5$ マスです。1マスが0.01なので、面積はおよそ
\[ 33.5\times0.01=0.335 \]
です。でも、黄のマスを「2つで1マス」と数えたのは、だいたいの決め方でした。マスの数え方で、答えが少し変わってしまいます。もっとはっきり決まる方法で、ほんとうの面積を追いつめていきましょう。

\begin{omoidasou}
\bkrecall{5-9}{$v$-$t$ グラフの面積は進んだ距離。方眼のマスや細い長方形の階段で考えた。}
\bkrecall{2-4}{分数のかけ算は、分子どうし・分母どうしをかける。$\left(\dfrac{2}{4}\right)^2=\dfrac{2}{4}\times\dfrac{2}{4}=\dfrac{4}{16}$}
\bkrecall{2-12}{$n^2\times n=n^3$（かける回数をたす）。まとめてわるときは、1つずつわってよい：$(A+B)\div h=A\div h+B\div h$}
\bkrecall{5-3}{$\displaystyle\lim_{n\to\infty}$ は「$n$ をかぎりなく大きくしたときの、近づく先」。$\infty$ は「無限大」。}
\bkrecall{4-5}{円の面積は、細いおうぎ形に切って並べかえ、長方形に近づけて求めた。}
\end{omoidasou}

%% ---------------- 数学の道具 ----------------
\bksection{doguu}{数学の道具}

\bksubsection{1}{長方形を4本ならべて、手で計算する}

\label{bk:5-10:hasami}マスのかわりに、たてに細長い長方形をならべます。まず、長方形を4本だけにして、1段ずつ考えましょう。

\noindent{\sffamily\bfseries\color{bkTheme}段1　はばを決める}\par
0から1までを、4つの等しいはばに分けます。長さ1を4つに等しく分けるので、1本のはばは
\[ 1\div4=\frac{1}{4} \]
です（わり算は分数で書ける、2-4）。

\noindent{\sffamily\bfseries\color{bkTheme}段2　$k$ 本めの右はしの場所を決める}\par
左から数えて1本め、2本め、3本め、4本めの長方形があります。この「何本め」を、文字 $k$ で表します。$k$ 本めの長方形の右はしは、0から、はば $\dfrac{1}{4}$ を $k$ こ進んだ所です。

\begin{bkfigure}
\begin{tikzpicture}[x=1mm, y=1mm]
  \draw[bkInk, line width=0.8pt] (0,0) -- (106,0);
  \foreach \k/\lab in {0/0, 1/\frac{1}{4}, 2/\frac{2}{4}, 3/\frac{3}{4}, 4/\frac{4}{4}} {
    \draw[bkInk, line width=0.8pt] (25*\k,-1.6) -- (25*\k,1.6);
    \node[font=\small, anchor=north, inner sep=1pt] at (25*\k,-2.6) {$\lab$};}
  \node[font=\sffamily\footnotesize, anchor=north, inner sep=1pt] at (100,-9.6) {（＝1）};
  \foreach \j in {1,2,3} {
    \draw[bkOrange!90!black, line width=0.9pt, -{Stealth[length=4pt]}] (25*\j-24,2.4) to[bend left=55] (25*\j-1,2.4);
    \node[font=\footnotesize, text=bkOrange!80!black, anchor=south, inner sep=1pt] at (25*\j-12.5,10.6) {$+\frac{1}{4}$};}
  \node[font=\sffamily\footnotesize, anchor=west, align=left, inner sep=0pt] at (78,9) {3本めの右はしは、\\$\frac{1}{4}$ を3こ進んだ $\frac{3}{4}$};
\end{tikzpicture}
\bkfigcaption{0から1までを4等分した目もり。はばを $k$ こ進むと、$k$ 本めの右はしに着く。}
\end{bkfigure}

だから、$k$ 本めの右はしの場所は
\[ x=\frac{1}{4}\times k=\frac{k}{4} \]
です。$k=3$ なら $\dfrac{3}{4}$、$k=4$ なら $\dfrac{4}{4}=1$ で、ちょうど右の終わりに着きます。

\noindent{\sffamily\bfseries\color{bkTheme}段3　長方形の高さを決める}\par
長方形の上の辺を、どの高さにするかを決めます。ここでは、長方形の\textbf{右上の角が、ちょうどグラフにのる}ようにします。これを\textbf{右はしの長方形}とよびます。右上の角は、右はし $x=\dfrac{k}{4}$ の真上にあるグラフの点です。グラフは $y=x^2$ なので、その点の高さは、$x^2$ の $x$ に $\dfrac{k}{4}$ を入れた
\[ y=\left(\frac{k}{4}\right)^2 \]
です。これが $k$ 本めの長方形の高さです。2乗がつくのは、グラフの式が $y=x^2$ だからです。

\begin{bkfigure}
\begin{tikzpicture}[x=84mm, y=42mm]
  \foreach \k in {1,...,4} {
    \fill[bkOrange!25] ({(\k-1)/4},0) rectangle ({\k/4},{(\k/4)^2});
    \draw[bkOrange!80!black, line width=0.5pt] ({(\k-1)/4},0) rectangle ({\k/4},{(\k/4)^2});
    \node[font=\small, anchor=south, inner sep=1pt] at ({(\k-0.5)/4},{(\k/4)^2+0.025}) {$\left(\frac{\k}{4}\right)^2$};
    \node[font=\sffamily\footnotesize, anchor=north, inner sep=1pt] at ({(\k-0.5)/4},-0.15) {\k 本め};}
  \draw[bkBlue!80!black, line width=1.3pt] plot[domain=0:1.04, samples=50] (\x,{\x*\x});
  \foreach \k in {1,...,4} {\fill[bkInk] ({\k/4},{(\k/4)^2}) circle (1.7pt);}
  \draw[bkInk, line width=0.7pt, -{Stealth[length=4pt]}] (0,0) -- (1.1,0) node[right, font=\sffamily\small, inner sep=1pt] {$x$};
  \draw[bkInk, line width=0.7pt, -{Stealth[length=4pt]}] (0,0) -- (0,1.16) node[above, font=\sffamily\small, inner sep=1pt] {$y$};
  \foreach \k in {1,...,4} {\node[font=\footnotesize, anchor=north, inner sep=1pt] at ({\k/4},-0.02) {$\frac{\k}{4}$};}
  \node[font=\sffamily\footnotesize, anchor=north east, inner sep=1pt] at (-0.005,-0.01) {0};
  \node[font=\sffamily\footnotesize, anchor=east, inner sep=1pt] at (-0.01,1) {1};
  \node[font=\sffamily\small, text=bkBlue!80!black, anchor=west, inner sep=1pt] at (1.05,1.05) {$y=x^2$};
\end{tikzpicture}
\bkfigcaption{右はしの長方形。黒い点（右上の角）がグラフにのる。長方形の上の数が高さ。}
\end{bkfigure}

\noindent{\sffamily\bfseries\color{bkTheme}段4　面積を計算して、たす}\par
長方形の面積は、高さ×はばです。分数のかけ算は、分子どうし・分母どうしをかけます（2-4）。表にまとめましょう。

\begin{bkcenter}
\begin{bktable}{colspec={Q[l,m,wd=34mm] *{4}{Q[c,m,wd=17mm]}}, rowsep=2pt}
$k$（何本め） & 1 & 2 & 3 & 4 \\
右はし $x=\dfrac{k}{4}$ & $\dfrac{1}{4}$ & $\dfrac{2}{4}$ & $\dfrac{3}{4}$ & $\dfrac{4}{4}$ \\
高さ $\left(\dfrac{k}{4}\right)^2$ & $\dfrac{1}{16}$ & $\dfrac{4}{16}$ & $\dfrac{9}{16}$ & $\dfrac{16}{16}$ \\
はば & $\dfrac{1}{4}$ & $\dfrac{1}{4}$ & $\dfrac{1}{4}$ & $\dfrac{1}{4}$ \\
面積＝高さ×はば & $\dfrac{1}{64}$ & $\dfrac{4}{64}$ & $\dfrac{9}{64}$ & $\dfrac{16}{64}$ \\
\end{bktable}
\bktabcaption{右はしの長方形4本。たとえば3本めは、$\frac{9}{16}\times\frac{1}{4}=\frac{9\times1}{16\times4}=\frac{9}{64}$}
\end{bkcenter}

4本の面積をたすと、分母がそろっているので、分子だけをたせます。
\[ \frac{1}{64}+\frac{4}{64}+\frac{9}{64}+\frac{16}{64}=\frac{30}{64}=0.46875 \]

\noindent{\sffamily\bfseries\color{bkTheme}左はしの長方形でも計算する}\par
こんどは、長方形の\textbf{左上の角}がグラフにのるようにします（\textbf{左はしの長方形}）。$k$ 本めの左はしは、右はしより1つ前の目もり $\dfrac{k-1}{4}$ なので、高さは $\left(\dfrac{0}{4}\right)^2$、$\left(\dfrac{1}{4}\right)^2$、$\left(\dfrac{2}{4}\right)^2$、$\left(\dfrac{3}{4}\right)^2$ です。1本めは高さが0なので、見えません。
\[ \frac{0}{64}+\frac{1}{64}+\frac{4}{64}+\frac{9}{64}=\frac{14}{64}=0.21875 \]

\begin{bkfigure}
\begin{tikzpicture}
\begin{axis}[bkplot, name=b1, width=60mm, height=48mm, xmin=0, xmax=1.15, ymin=0, ymax=1.15, xtick={0.25,0.5,0.75,1}, ytick={0.5,1}, xlabel={$x$}, ylabel={$y$}, title={\sffamily\small 右はしの長方形（大きめ）}]
  \foreach \k in {1,...,4} {\edef\tmp{\noexpand\fill[bkOrange!30, draw=bkOrange!70, line width=0.5pt] (axis cs:{(\k-1)/4},0) rectangle (axis cs:{\k/4},{(\k/4)^2});}\tmp}
  \addplot[bkcurve, domain=0:1.05] {x^2};
\end{axis}
\begin{axis}[bkplot, at={(b1.east)}, anchor=west, xshift=14mm, width=60mm, height=48mm, xmin=0, xmax=1.15, ymin=0, ymax=1.15, xtick={0.25,0.5,0.75,1}, ytick={0.5,1}, xlabel={$x$}, ylabel={$y$}, title={\sffamily\small 左はしの長方形（小さめ）}]
  \foreach \k in {1,...,4} {\edef\tmp{\noexpand\fill[bkBlue!25, draw=bkBlue!70, line width=0.5pt] (axis cs:{(\k-1)/4},0) rectangle (axis cs:{\k/4},{((\k-1)/4)^2});}\tmp}
  \addplot[bkcurve, domain=0:1.05] {x^2};
\end{axis}
\end{tikzpicture}
\bkfigcaption{$n=4$ の場合。右はしの長方形はグラフの上にはみ出し、左はしの長方形はグラフの下にすき間がある。}
\end{bkfigure}

右はしの長方形は、グラフの上に少しはみ出しています。だから、ほんとうの面積より大きい。左はしの長方形は、グラフの下にすき間があります。だから、ほんとうの面積より小さい。ほんとうの面積を $S$ とすると、
\[ 0.21875<S<0.46875 \]
です。$S$ は、2つの数のあいだに\textbf{はさまれて}います。やってみようの0.335も、このあいだにあります。

\bksubsection{2}{長方形をふやして、はさみうちにする}

4本では、はさんだはばが $0.46875-0.21875=0.25$ もあって、ゆるすぎます。長方形を細くして、本数をふやしましょう。本数を $n$ 本にしても、段1〜段4は同じです。

\begin{bkcenter}
\begin{bktable}{colspec={Q[l,m,wd=30mm] Q[c,m,wd=26mm] Q[c,m,wd=26mm] Q[c,m,wd=34mm]}, rowsep=2pt, row{1}={font=\sffamily\bfseries\small}}
 & 4本 & 10本 & $n$ 本 \\
はば & $\dfrac{1}{4}$ & $\dfrac{1}{10}$ & $\dfrac{1}{n}$\quad（$1\div n$） \\
$k$ 本めの右はし & $\dfrac{k}{4}$ & $\dfrac{k}{10}$ & $\dfrac{k}{n}$ \\
$k$ 本めの高さ & $\left(\dfrac{k}{4}\right)^2$ & $\left(\dfrac{k}{10}\right)^2$ & $\left(\dfrac{k}{n}\right)^2$ \\
$k$ 本めの面積 & $\left(\dfrac{k}{4}\right)^2\times\dfrac{1}{4}$ & $\left(\dfrac{k}{10}\right)^2\times\dfrac{1}{10}$ & $\left(\dfrac{k}{n}\right)^2\times\dfrac{1}{n}$ \\
\end{bktable}
\bktabcaption{本数が4でも10でも $n$ でも、考え方は同じ。4 のところが $n$ に変わるだけ。}
\end{bkcenter}

10本なら、表計算（2-18）で計算できます。A列に $k$、B列に右はし、C列に高さ、D列に面積を入れます。

\begin{bkcenter}
\bksheet{4}{24}{%
1 & $k$ & 右はし $x$ & 高さ $x^2$ & 面積 \\
2 & 1 & 0.1 & 0.01 & 0.001 \\
3 & 2 & 0.2 & 0.04 & 0.004 \\
4 & 3 & 0.3 & 0.09 & 0.009 \\
5 & 4 & 0.4 & 0.16 & 0.016 \\
6 & 5 & 0.5 & 0.25 & 0.025 \\
7 & 6 & 0.6 & 0.36 & 0.036 \\
8 & 7 & 0.7 & 0.49 & 0.049 \\
9 & 8 & 0.8 & 0.64 & 0.064 \\
10 & 9 & 0.9 & 0.81 & 0.081 \\
11 & 10 & 1 & 1 & 0.1 \\
12 & & & 合計 & 0.385 \\}
\bktabcaption{B2 に \texttt{=A2/10}（右はし $\frac{k}{10}$）、C2 に \texttt{=B2\^{}2}（2乗）、D2 に \texttt{=C2*0.1}（はば0.1をかける）と入れて、下へコピーする。}
\end{bkcenter}

D12 には \texttt{=SUM(D2:D11)} と入れます。SUM は英語で「和（たし算の答え）」という意味で、「D2 から D11 までを全部たす」という命令です。右はしの和は0.385になりました。左はしの和は、同じように計算して0.285です。

$n$ を100、1000とふやすと、次の表のようになります（「2つの平均」は、長方形のかわりに台形をならべた和です。あとの例題8で使います）。

\begin{bkcenter}
\bksheet{4}{24}{%
1 & $n$ & 左はしの和 & 右はしの和 & 2つの平均 \\
2 & 4 & 0.21875 & 0.46875 & 0.34375 \\
3 & 10 & 0.285 & 0.385 & 0.335 \\
4 & 100 & 0.32835 & 0.33835 & 0.33335 \\
5 & 1000 & 0.3328335 & 0.3338335 & 0.3333335 \\}
\bktabcaption{$n$ を大きくすると、左はしの和と右はしの和が、両がわから同じ数に近づいていく。}
\end{bkcenter}

\noindent{\sffamily\bfseries\color{bkTheme}2つの和の差は、いつも $\dfrac{1}{n}$}\par
表を見ると、右はしの和と左はしの和の差は、$n=4$ で0.25、$n=10$ で0.1、$n=100$ で0.01です。いつも $\dfrac{1}{n}$ になっています。そのわけを、図で見ましょう。

\begin{bkfigure}
\begin{tikzpicture}[x=40mm, y=34mm]
  \begin{scope}
    \fill[bkGreen!45] (0.25,0) rectangle (0.5,0.0625);
    \fill[bkBlue!35] (0.5,0) rectangle (0.75,0.25);
    \fill[bkPurple!35] (0.75,0) rectangle (1,0.5625);
    \draw[bkInk!70, line width=0.4pt] (0.25,0) rectangle (0.5,0.0625) (0.5,0) rectangle (0.75,0.25) (0.75,0) rectangle (1,0.5625);
    \draw[bkBlue!80!black, line width=1.1pt] plot[domain=0:1.03, samples=40] (\x,{\x*\x});
    \draw[bkInk, line width=0.6pt] (0,0) -- (1.06,0) (0,0) -- (0,1.08);
    \node[font=\sffamily\small, anchor=south, inner sep=1pt] at (0.5,1.1) {左はしの長方形};
    \node[font=\sffamily\footnotesize, anchor=north, inner sep=1pt] at (0,-0.02) {0};
    \node[font=\sffamily\footnotesize, anchor=north, inner sep=1pt] at (1,-0.02) {1};
  \end{scope}
  \draw[bk arrow, bkGray, line width=0.9pt] (1.12,0.4) -- (1.36,0.4);
  \node[font=\sffamily\footnotesize, anchor=south, align=center, inner sep=1pt] at (1.24,0.47) {1つ左へ\\ずらす};
  \begin{scope}[shift={(1.5,0)}]
    \fill[bkGreen!45] (0,0) rectangle (0.25,0.0625);
    \fill[bkBlue!35] (0.25,0) rectangle (0.5,0.25);
    \fill[bkPurple!35] (0.5,0) rectangle (0.75,0.5625);
    \fill[bkRed!50] (0.75,0) rectangle (1,1);
    \draw[bkInk!70, line width=0.4pt] (0,0) rectangle (0.25,0.0625) (0.25,0) rectangle (0.5,0.25) (0.5,0) rectangle (0.75,0.5625) (0.75,0) rectangle (1,1);
    \draw[bkBlue!80!black, line width=1.1pt] plot[domain=0:1.03, samples=40] (\x,{\x*\x});
    \draw[bkInk, line width=0.6pt] (0,0) -- (1.06,0) (0,0) -- (0,1.08);
    \node[font=\sffamily\small, anchor=south, inner sep=1pt] at (0.5,1.1) {右はしの長方形};
    \node[font=\sffamily\footnotesize, anchor=north, inner sep=1pt] at (0,-0.02) {0};
    \node[font=\sffamily\footnotesize, anchor=north, inner sep=1pt] at (1,-0.02) {1};
    \node[font=\sffamily\footnotesize, anchor=west, align=left, inner sep=0pt] at (1.1,0.55) {赤の1本だけ\\多い。\\はば $\frac{1}{4}$、\\高さ1};
  \end{scope}
\end{tikzpicture}
\bkfigcaption{同じ色の長方形は、同じ形。右はしの長方形は、左はしの長方形を1つ左へずらして、赤の1本を足したもの。}
\end{bkfigure}

左はしの長方形の2本め（緑）は、右はしの長方形の1本めと同じ高さ $\left(\frac{1}{4}\right)^2$ です。3本め（青）と2本め、4本め（むらさき）と3本めも同じです。だから、ちがうのは右はしの4本め（赤）だけです。赤の長方形は、高さ $1^2=1$、はば $\frac{1}{n}$ なので、面積は $1\times\dfrac{1}{n}=\dfrac{1}{n}$ です。

$n=1000$ なら、差はたったの0.001です。ほんとうの面積 $S$ は
\[ 0.3328335<S<0.3338335 \]
のあいだにはさまれています。$n$ をもっと大きくすると、差 $\dfrac{1}{n}$ はいくらでも小さくなり、2つの和は同じ数に近づきます。表の数から、近づく先は $\dfrac{1}{3}=0.3333\cdots$ のようです。ほんとうに $\dfrac{1}{3}$ なのかを、④で式を使って確かめます。

\bksubsection{3}{たし算をまとめて書く：$\Sigma$}

\label{bk:5-10:sigma}長方形が1000本あると、たし算を全部書くのは大変です。そこで、「番号を変えながら、全部たす」ことを短く書く記号を使います。

\begin{bkfigure}
\begin{tikzpicture}[x=1mm, y=1mm]
  \node[inner sep=0pt, font=\LARGE] (S) at (0,0) {$\displaystyle\sum$};
  \node[inner sep=1pt, anchor=south, font=\normalsize] (up) at ([yshift=0.8mm]S.north) {$4$};
  \node[inner sep=1pt, anchor=north, font=\normalsize] (lo) at ([yshift=-0.8mm]S.south) {$k=1$};
  \node[inner sep=1pt, anchor=west, font=\LARGE] (b) at ([xshift=1mm]S.east) {$k^2$};
  \tikzset{lab/.style={font=\sffamily\footnotesize, text=bkInk, inner sep=1pt, align=left}, ar/.style={bkTheme, line width=0.7pt, -{Stealth[length=4pt]}, shorten >=1.2mm}}
  \node[lab, anchor=east] (l1) at (-16,9) {おわりの番号（4まで）};
  \node[lab, anchor=east] (l2) at (-16,-9) {はじめの番号（$k=1$ から）};
  \node[lab, anchor=west] (l3) at (20,0) {たすもの\\（$k$ に番号を入れる）};
  \draw[ar] (l1.east) -- (up.west);
  \draw[ar] (l2.east) -- (lo.west);
  \draw[ar] (l3.west) -- (b.east);
\end{tikzpicture}
\bkfigcaption{$\Sigma$ の読み方。「$k$ を1から4まで1ずつ変えて、$k^2$ を全部たす」}
\end{bkfigure}

この記号は、次のように計算します。

\begin{bkcenter}
\begin{bktable}{colspec={Q[c,m,wd=22mm] Q[c,m,wd=26mm] Q[c,m,wd=20mm]}, rowsep=1.5pt, row{1}={font=\sffamily\bfseries\small}}
$k$ に入れる数 & $k^2$ & 値 \\
1 & $1^2$ & 1 \\
2 & $2^2$ & 4 \\
3 & $3^2$ & 9 \\
4 & $4^2$ & 16 \\
\SetCell[c=2]{r} 全部たす & & 30 \\
\end{bktable}
\end{bkcenter}
\[ \sum_{k=1}^{4}k^2=1^2+2^2+3^2+4^2=1+4+9+16=30 \]

$\Sigma$ は、ギリシャ文字の「シグマ」で、英語の sum（和）の頭文字 S にあたります。表計算の SUM と同じ意味です。上の式は「シグマ、$k$ イコール1から4まで、$k$ の2乗」と読みます。\bkidx{しぐま}{$\Sigma$（シグマ）}

\noindent\textbf{例1}\quad $\displaystyle\sum_{k=1}^{3}k=1+2+3=6$\quad（$k$ そのものをたす）\par
\noindent\textbf{例2}\quad $\displaystyle\sum_{k=1}^{4}(2k+1)=3+5+7+9=24$\quad（$k=1,2,3,4$ を順に入れる。$k=1$ なら $2\times1+1=3$）

\noindent{\sffamily\bfseries\color{bkTheme}長方形の和を $\Sigma$ で書く}\par
4本の右はしの長方形の和は、$k$ 本めの面積 $\left(\dfrac{k}{4}\right)^2\times\dfrac{1}{4}$ を、$k=1$ から4までたしたものです。
\[ \sum_{k=1}^{4}\left(\frac{k}{4}\right)^2\times\frac{1}{4}=\left(\frac{1}{4}\right)^2\times\frac{1}{4}+\left(\frac{2}{4}\right)^2\times\frac{1}{4}+\left(\frac{3}{4}\right)^2\times\frac{1}{4}+\left(\frac{4}{4}\right)^2\times\frac{1}{4} \]
$n$ 本なら、4を $n$ にかえて、
\[ S_n=\sum_{k=1}^{n}\left(\frac{k}{n}\right)^2\times\frac{1}{n} \]
と書けます。$S_n$ は「$n$ 本のときの、右はしの長方形の和」という名前です。

\noindent{\sffamily\bfseries\color{bkTheme}1から $n$ までの和}\par
あとで使うので、$\displaystyle\sum_{k=1}^{n}k=1+2+\cdots+n$ を計算しておきます。$n=4$ で、ブロックの階段を考えます。

\begin{bkfigure}
\begin{tikzpicture}[x=6mm, y=6mm]
  \foreach \j in {1,...,4} {\foreach \i in {1,...,\j} {{\draw[bkBlue!80!black, fill=bkBlue!30, line width=0.5pt] (\j-1,\i-1) rectangle (\j,\i);}}}
  \foreach \j in {1,...,4} {\node[font=\sffamily\footnotesize, anchor=north, inner sep=1pt] at (\j-0.5,-0.15) {\j};}
  \node[font=\sffamily\small, anchor=north, inner sep=1pt] at (2,-1.1) {$1+2+3+4$};
  \draw[bk arrow, bkGray, line width=0.9pt] (5,2.5) -- (7,2.5);
  \node[font=\sffamily\footnotesize, anchor=south, align=center, inner sep=1pt] at (6,2.8) {同じ階段を\\さかさにのせる};
  \begin{scope}[shift={(8,0)}]
    \foreach \j in {1,...,4} {\foreach \i in {1,...,\j} {{\draw[bkBlue!80!black, fill=bkBlue!30, line width=0.5pt] (\j-1,\i-1) rectangle (\j,\i);}}}
    \foreach \j in {1,...,4} {\foreach \i in {\j,...,4} {{\draw[bkOrange!80!black, fill=bkOrange!35, line width=0.5pt] (\j-1,\i) rectangle (\j,\i+1);}}}
    \draw[bkGray, line width=0.5pt, <->] (0,-0.5) -- (4,-0.5);
    \node[font=\sffamily\footnotesize, anchor=north, inner sep=1pt] at (2,-0.7) {よこ4};
    \draw[bkGray, line width=0.5pt, <->] (4.5,0) -- (4.5,5);
    \node[font=\sffamily\footnotesize, anchor=west, inner sep=1pt] at (4.7,2.5) {たて5（$4+1$）};
  \end{scope}
\end{tikzpicture}
\bkfigcaption{階段を2つ組み合わせると、よこ4・たて5の長方形になる。}
\end{bkfigure}

長方形のブロックは $4\times5=20$ こ。階段はその半分なので、
\[ 1+2+3+4=20\div2=10 \]
です。$n$ 段の階段でも、2つ組み合わせると、よこ $n$、たて $n+1$ の長方形になります。だから
\[ \sum_{k=1}^{n}k=1+2+\cdots+n=\frac{n(n+1)}{2} \]
です。$n=3$ なら $\dfrac{3\times4}{2}=6$ で、例1と合います。

\bksubsection{4}{$n$ を大きくした先を、式で求める}

\label{bk:5-10:kyokugen}②の表では、右はしの和 $S_n$ は $\dfrac{1}{3}$ に近づくようでした。$S_n$ を $n$ の式で表して、ほんとうに $\dfrac{1}{3}$ になるかを確かめます。

\noindent{\sffamily\bfseries\color{bkTheme}段1　1本の面積を、1つの分数にする}\par
分数の2乗は、分子と分母をそれぞれ2乗します（2-4）。そのあと $\dfrac{1}{n}$ をかけると、分母は $n^2\times n=n^3$ になります（2-12）。
\[ \left(\frac{k}{n}\right)^2\times\frac{1}{n}=\frac{k^2}{n^2}\times\frac{1}{n}=\frac{k^2}{n^3} \]
$n=4$、$k=2$ で確かめると、$\dfrac{2^2}{4^3}=\dfrac{4}{64}$。①の表の2本めと合います。

\noindent{\sffamily\bfseries\color{bkTheme}段2　分母が同じなので、分子だけをたす}\par
どの長方形も分母が $n^3$ でそろっているので、分子だけをたせます。①で $\dfrac{1+4+9+16}{64}$ としたのと同じです。
\[ S_n=\frac{1^2}{n^3}+\frac{2^2}{n^3}+\cdots+\frac{n^2}{n^3}=\frac{1^2+2^2+\cdots+n^2}{n^3} \]

\noindent{\sffamily\bfseries\color{bkTheme}段3　2乗の和を、$n$ の式にする}\par
分子の $1^2+2^2+\cdots+n^2$ には、次の式があります。まず、表で合っているかを確かめましょう。
\[ 1^2+2^2+\cdots+n^2=\frac{n(n+1)(2n+1)}{6} \]

\begin{bkcenter}
\begin{bktable}{colspec={Q[l,m,wd=44mm] *{6}{Q[c,m,wd=10mm]}}, rowsep=1.5pt}
$n$ & 1 & 2 & 3 & 4 & 5 & 6 \\
$1^2+2^2+\cdots+n^2$ & 1 & 5 & 14 & 30 & 55 & 91 \\
$\dfrac{n(n+1)(2n+1)}{6}$ & 1 & 5 & 14 & 30 & 55 & 91 \\
\end{bktable}
\bktabcaption{たとえば $n=4$ なら $\frac{4\times5\times9}{6}=\frac{180}{6}=30$。どの $n$ でも合う。}
\end{bkcenter}

なぜこの式になるのかは、もっと知りたいで説明します。

\begin{motto}[4]
\noindent\textbf{2乗の和の式のわけ}\par
5-4の立方体の図で、1辺 $x$ の立方体を $h$ だけ大きくすると、$(x+h)^3=x^3+3x^2h+3xh^2+h^3$ でした。$x=k$、$h=1$ とすると、
\[ (k+1)^3-k^3=3k^2+3k+1 \]
です。$k=1,\,2,\,3,\,4$ を入れて、たてにたしてみます。
\begin{align*}
2^3-1^3&=3\times1^2+3\times1+1\\
3^3-2^3&=3\times2^2+3\times2+1\\
4^3-3^3&=3\times3^2+3\times3+1\\
5^3-4^3&=3\times4^2+3\times4+1
\end{align*}
左がわをたすと、$2^3$ と $-2^3$、$3^3$ と $-3^3$、$4^3$ と $-4^3$ が消えて、$5^3-1^3=124$ だけが残ります。右がわは、次の式です。
\[ 3\times(1^2+2^2+3^2+4^2)+3\times(1+2+3+4)+4 \]
\noindent $1+2+3+4=10$（③）なので、
\[ 124=3\times(1^2+2^2+3^2+4^2)+30+4 \]
これから $1^2+2^2+3^2+4^2=(124-34)\div3=30$。4を $n$ にかえて同じ計算をすると、$(n+1)^3-1=3\times(1^2+\cdots+n^2)+3\times\dfrac{n(n+1)}{2}+n$ となり、これを整理すると上の式になります。
\end{motto}

\noindent{\sffamily\bfseries\color{bkTheme}段4　$S_n$ を、$\dfrac{1}{n}$ の式に直す}\par
段2と段3から、
\[ S_n=\frac{n(n+1)(2n+1)}{6}\div n^3=\frac{n(n+1)(2n+1)}{6\times n\times n\times n} \]
分母の $n\times n\times n$ の $n$ を1つずつ、分子のかっこに配ります。
\[ S_n=\frac{1}{6}\times\frac{n}{n}\times\frac{n+1}{n}\times\frac{2n+1}{n} \]
ここで、$\dfrac{n}{n}=1$。そして、まとめてわるときは1つずつわってよい（2-12）ので、
\[ \frac{n+1}{n}=\frac{n}{n}+\frac{1}{n}=1+\frac{1}{n},\qquad \frac{2n+1}{n}=\frac{2n}{n}+\frac{1}{n}=2+\frac{1}{n} \]
です。だから
\[ S_n=\frac{1}{6}\Bigl(1+\frac{1}{n}\Bigr)\Bigl(2+\frac{1}{n}\Bigr) \]
となります。$n=4$ を入れると、$\dfrac{1}{6}\times1.25\times2.25=0.46875$。$n=10$ なら $\dfrac{1}{6}\times1.1\times2.1=0.385$。②の表とぴったり合います。

\noindent{\sffamily\bfseries\color{bkTheme}段5　$n$ をかぎりなく大きくする}\par
$n$ を大きくすると、$\dfrac{1}{n}$ は0.25、0.1、0.01、0.001、……と、0にかぎりなく近づきます（5-3）。だから、$1+\dfrac{1}{n}$ は1に、$2+\dfrac{1}{n}$ は2に近づきます。
\[ \lim_{n\to\infty}S_n=\frac{1}{6}\times1\times2=\frac{2}{6}=\frac{1}{3} \]
表計算で見つけた $\dfrac{1}{3}$ が、式でも確かめられました。左はしの和は、右はしの和より $\dfrac{1}{n}$ 小さいだけ（②）なので、$\dfrac{1}{3}-0=\dfrac{1}{3}$ に近づきます。左右からはさまれたほんとうの面積 $S$ は、ちょうど $\dfrac{1}{3}$ です。

\bkhatenaref の答え：$v=t^2$ のグラフの、0秒から1秒までの面積は $\dfrac{1}{3}$ なので、ボールは0秒から1秒までに $\dfrac{1}{3}\,\text{m}$（約\qty{0.33}{m}）進みます。やってみようの0.335とも、よく合っています。

\bksubsection{5}{はばを0に近づけた先が、定積分}

\label{bk:5-10:teisekibun}はばをせまくしていくようすを、図で見てみましょう。

\begin{bkfigure}
\begin{tikzpicture}[x=1mm, y=1mm]
  \foreach \n/\px/\w/\sm/\ov in {4/0/{$\frac{1}{4}$}/0.469/0.135, 8/36/{$\frac{1}{8}$}/0.398/0.065, 16/72/{$\frac{1}{16}$}/0.365/0.032} {
    \begin{scope}[shift={(\px,8)}]
      % 長方形（右はし）。はみ出しは こい色＋斜線
      \foreach \k in {1,...,\n} {
        \fill[bkOrange!70] ({26*(\k-1)/\n},0) rectangle ({26*\k/\n},{22*(\k/\n)^2});
        \fill[pattern=north east lines, pattern color=white] ({26*(\k-1)/\n},0) rectangle ({26*\k/\n},{22*(\k/\n)^2});}
      \fill[bkOrange!25] (0,0) -- plot[domain=0:1, samples=40] ({26*\x},{22*\x*\x}) -- (26,0) -- cycle;
      \foreach \k in {1,...,\n} {\draw[bkOrange!80!black, line width=0.3pt] ({26*(\k-1)/\n},0) rectangle ({26*\k/\n},{22*(\k/\n)^2});}
      \draw[bkBlue!80!black, line width=1pt] plot[domain=0:1.04, samples=40] ({26*\x},{22*\x*\x});
      \draw[bkInk, line width=0.6pt] (0,0) -- (29,0) (0,0) -- (0,25);
      \node[font=\sffamily\footnotesize, anchor=north, inner sep=1pt] at (26,-0.3) {1};
      \node[font=\sffamily\bfseries\footnotesize, anchor=north west, inner sep=0pt, align=left] at (0,-4.2) {はば \w};
      \node[font=\sffamily\footnotesize, anchor=north west, inner sep=0pt, align=left] at (0,-10) {和 \sm\\{\color{bkOrange!80!black}はみ出し \ov}};
    \end{scope}
    \draw[bk arrow, bkGray, line width=0.8pt] (\px+28.5,19) -- (\px+34.5,19);}
  % はばを0に近づけた先
  \begin{scope}[shift={(108,8)}]
    \fill[bkYellow!40] (0,0) -- plot[domain=0:1, samples=40] ({26*\x},{22*\x*\x}) -- (26,0) -- cycle;
    \draw[bkBlue!80!black, line width=1pt] plot[domain=0:1.04, samples=40] ({26*\x},{22*\x*\x});
    \draw[bkInk, line width=0.6pt] (0,0) -- (29,0) (0,0) -- (0,25);
    \node[font=\sffamily\footnotesize, anchor=north, inner sep=1pt] at (26,-0.3) {1};
    \node[font=\sffamily\bfseries\small] at (17.5,4.5) {$\frac{1}{3}$};
    \node[font=\sffamily\bfseries\footnotesize, anchor=north west, inner sep=0pt, align=left] at (0,-4.2) {はば → 0};
    \node[font=\sffamily\footnotesize, anchor=north west, inner sep=0pt, align=left] at (0,-10) {和 → $\frac{1}{3}$\\{\color{bkOrange!80!black}はみ出し → 0}};
  \end{scope}
\end{tikzpicture}
\bkfigcaption{はばを半分にするたびに、はみ出し（こい色の斜線）は、ほぼ半分になる。はばを0に近づけた先は、曲線の下の面積そのもの。}
\end{bkfigure}

はばを0に近づけていくと、階段の面積は、はてな？の色をぬった部分の面積そのものになります。ここまでにしたことは、いつも次の3段です。

\begin{bkfigure}
\begin{tikzpicture}[x=1mm, y=1mm]
  \tikzset{cv/.style={bkBlue!80!black, line width=1pt}, ax/.style={bkInk, line width=0.6pt}}
  \def\fx(#1){(6+14*(#1)*(#1))}
  \foreach \px/\ttl in {0/{① 細かく分ける}, 46/{② 1本の面積をかける}, 92/{③ 全部たす}} {
    \node[font=\sffamily\bfseries\small, anchor=south, inner sep=1pt] at (\px+16,25) {\ttl};}
  % ① 分ける
  \begin{scope}[shift={(0,0)}]
    \fill[bkYellow!35] (2,0) -- plot[domain=0.0667:1, samples=30] ({30*\x},{6+14*\x*\x}) -- (30,0) -- cycle;
    \foreach \k in {1,...,6} {\draw[bkGray, line width=0.4pt, dash pattern=on 1.2pt off 1pt] ({2+28*\k/6},0) -- ({2+28*\k/6},{6+14*((2+28*\k/6)/30)^2});}
    \draw[cv] plot[domain=0.0667:1.03, samples=30] ({30*\x},{6+14*\x*\x});
    \draw[ax] (0,0) -- (33,0) (0,0) -- (0,23);
    \draw[bkGray, line width=0.5pt, <->] (11.33,-1.6) -- (16,-1.6);
    \node[font=\footnotesize, anchor=north, inner sep=1pt] at (13.67,-2.2) {$\Delta x$};
    \node[font=\footnotesize, anchor=north, inner sep=1pt] at (2,-2.2) {$a$};
    \node[font=\footnotesize, anchor=north, inner sep=1pt] at (30,-2.2) {$b$};
  \end{scope}
  \draw[bk arrow, bkGray, line width=0.8pt] (35,11) -- (42,11);
  % ② かける
  \begin{scope}[shift={(46,0)}]
    \fill[bkYellow!35] (2,0) -- plot[domain=0.0667:1, samples=30] ({30*\x},{6+14*\x*\x}) -- (30,0) -- cycle;
    \fill[bkOrange!45] (16,0) rectangle (20.67,{6+14*(20.67/30)^2});
    \draw[bkOrange!80!black, line width=0.5pt] (16,0) rectangle (20.67,{6+14*(20.67/30)^2});
    \draw[cv] plot[domain=0.0667:1.03, samples=30] ({30*\x},{6+14*\x*\x});
    \draw[ax] (0,0) -- (33,0) (0,0) -- (0,23);
    \draw[bkGray, line width=0.5pt, <->] (16,-1.6) -- (20.67,-1.6);
    \node[font=\footnotesize, anchor=north, inner sep=1pt] at (18.33,-2.2) {$\Delta x$};
    \draw[bkGray, line width=0.5pt, <->] (14.2,0) -- (14.2,{6+14*(20.67/30)^2});
    \node[font=\footnotesize, anchor=east, inner sep=1pt, fill=bkYellow!35] at (13.6,3.2) {$f(x)$};
    \node[font=\sffamily\footnotesize, anchor=north, inner sep=1pt] at (16,-7.5) {高さ×はば＝$f(x)\,\Delta x$};
  \end{scope}
  \draw[bk arrow, bkGray, line width=0.8pt] (81,11) -- (88,11);
  % ③ たす
  \begin{scope}[shift={(92,0)}]
    \foreach \k in {1,...,6} {\fill[bkOrange!35] ({2+28*(\k-1)/6},0) rectangle ({2+28*\k/6},{6+14*((2+28*\k/6)/30)^2});
      \draw[bkOrange!80!black, line width=0.4pt] ({2+28*(\k-1)/6},0) rectangle ({2+28*\k/6},{6+14*((2+28*\k/6)/30)^2});}
    \draw[cv] plot[domain=0.0667:1.03, samples=30] ({30*\x},{6+14*\x*\x});
    \draw[ax] (0,0) -- (33,0) (0,0) -- (0,23);
    \node[font=\sffamily\footnotesize, anchor=north, align=center, inner sep=1pt] at (16,-2.2) {$\sum f(x)\,\Delta x$\\はばを0に近づけて\\$\displaystyle\int_a^b f(x)\,\dd x$};
  \end{scope}
\end{tikzpicture}
\bkfigcaption{積分の3段。① はば $\Delta x$ に細かく分ける。② 1本ずつ「高さ×はば」をかける。③ 全部たして、はばを0に近づける。}
\end{bkfigure}

この近づく先を、次のように書きます。

\begin{bkfigure}
\begin{tikzpicture}[x=1mm, y=1mm]
  \node[inner sep=0pt, font=\Large] (i) at (0,0) {$\displaystyle\int_0^1$};
  \node[inner sep=0pt, font=\Large, anchor=west] (f) at ([xshift=0.6mm]i.east) {$x^2$};
  \node[inner sep=0pt, font=\Large, anchor=west] (d) at ([xshift=1.4mm]f.east) {$\dd x$};
  \node[inner sep=0pt, font=\Large, anchor=west] (e) at ([xshift=1.4mm]d.east) {$=\dfrac{1}{3}$};
  \tikzset{lab/.style={font=\sffamily\footnotesize, text=bkInk, inner sep=1pt, align=left}, ar/.style={bkTheme, line width=0.7pt, -{Stealth[length=4pt]}, shorten >=1.2mm}}
  % 左がわ：∫ と 0・1
  \node[lab, anchor=east] (l1) at ([xshift=-12mm, yshift=6mm]i.center) {$\int$ ＝ たし合わせる\\（sum の S をのばした形）};
  \node[lab, anchor=east] (l2) at ([xshift=-12mm, yshift=-7mm]i.center) {下の 0、上の 1 ＝\\$x=0$ から 1 まで};
  \draw[ar] (l1.east) -- ([xshift=-0.8mm, yshift=2mm]i.center);
  \draw[ar] (l2.east) -- ([xshift=-2mm, yshift=-4.5mm]i.center);
  % 右がわ：式の右に列をつくる。上の2つは式の上を通って下ろす（式に線を重ねない）
  \coordinate (R) at ([xshift=9mm]e.east);
  \node[lab, anchor=west] (l3) at ([yshift=14mm]R |- i.center) {$x^2$ ＝ 長方形の高さ};
  \node[lab, anchor=west] (l4) at ([yshift=9mm]R |- i.center) {$\dd x$ ＝ かぎりなく細いはば};
  \node[lab, anchor=west] (l5) at ([yshift=-1mm]R |- i.center) {$\frac{1}{3}$ ＝ ④で求めた、近づく先};
  \draw[ar] (l3.west) -| (f.north);
  \draw[ar] (l4.west) -| (d.north);
  \draw[ar] (l5.west) -- (e.east |- l5.west);
\end{tikzpicture}
\bkfigcaption{定積分の記号の読み方}
\end{bkfigure}

$\int$ は「インテグラル」と読みます。上の式は「インテグラル、0から1まで、$x$ の2乗、ディーエックス」と読みます。

一般に、グラフ $y=f(x)$ の $x=a$ から $b$ までを細い長方形に分けます。$k$ 本めの長方形の右はしを $x_k$、はばを $\Delta x$ と書きます（④では $x_k=\dfrac{k}{n}$、$\Delta x=\dfrac{1}{n}$。たとえば $n=4$ なら $x_1=\dfrac{1}{4}$、$x_2=\dfrac{2}{4}$、$x_3=\dfrac{3}{4}$、$x_4=\dfrac{4}{4}=1$ で、はばはどれも $\Delta x=\dfrac{1}{4}$）。$\Delta x$ は「$x$ の変化量」（2-3）で、長方形の左はしから右はしまで、$x$ がどれだけ変わったかです。長方形の和は $\sum f(x_k)\,\Delta x$ です。このはばを0に近づけた先を
\[ \sum f(x_k)\,\Delta x\quad\longrightarrow\quad\int_a^b f(x)\,\dd x \]
と書き、$f(x)$ の $a$ から $b$ までの\textbf{定積分}（ていせきぶん）といいます。$\Sigma$ が $\int$ に、はば $\Delta x$ が $\dd x$ に変わったと考えましょう。$a$ を\textbf{下端}、$b$ を\textbf{上端}といいます。\bkidx{ていせきぶん}{定積分}\bkidx{いんてぐらる}{$\int$（インテグラル）}

\noindent{\sffamily\bfseries\color{bkTheme}なぜ $\dd x$ を書くのか}\par
\begin{enumerate}[label={\bkcircled[bkTheme]{\arabic*}}, leftmargin=2.4\zw, itemsep=1pt, topsep=2pt]
\item $f(x)\,\dd x$ は「高さ×はば」、つまり細い長方形1本の面積です。$\dd x$ がないと、たすものが面積になりません。
\item 何を細かく分けたのかを表します。$\displaystyle\int v\,\dd t$ なら、時間 $t$ を細かく分けています。
\item 単位がわかります。$v$（m/s）に $\dd t$（s）をかけるので、$(\text{m/s})\times\text{s}=\text{m}$。答えは距離になります（5-9）。
\end{enumerate}

\noindent{\sffamily\bfseries\color{bkTheme}マイナスになるとき}\par
$f(x)$ がマイナスの所では、長方形の高さがマイナスなので、その面積もマイナスとして数えます（5-9のマイナスの面積と同じ）。

また、下端と上端を入れかえて $\displaystyle\int_1^0 x^2\,\dd x$ と書くと、1から0へ、左向きに分けていくことになります。はば $\Delta x$ は「あと${}-{}$まえ」（2-3）なので、4本に分けるなら $\dfrac{0-1}{4}=-\dfrac{1}{4}$ で、マイナスです。高さ×はばがどれもマイナスになるので、
\[ \int_1^0 x^2\,\dd x=-\frac{1}{3} \]
です。下端と上端を入れかえると、符号が反対になります（5-11でもう一度確かめます）。

\begin{seikaku}
\noindent\textbf{はばと高さの決め方}\par
ここでは、はばをすべて同じにし、高さに右はし（または左はし）の値を使いました。せいかくには、はばはそろっていなくてもよく、高さには小さな区間の中のどこの値を使ってもかまいません。グラフがつながっている（とぎれていない）関数なら、いちばん広いはばを0に近づけたとき、どの決め方でも和は同じ値に近づきます。それが定積分です。
\end{seikaku}

\bkkigocards{\bkkigocard{$\displaystyle\sum_{k=1}^{n}$}{シグマ}{$k$ を1から $n$ まで変えて、全部たす}{sum}{5-10}}
{\bkkigocard{$\displaystyle\int_a^b$}{インテグラル}{$a$ から $b$ まで、細かく分けてたす（定積分）}{integral}{5-10}}

\bksubsection{6}{0から $b$ までの面積}

\label{bk:5-10:b}④では、0から1までを求めました。0から2までなら、どうなるでしょう。同じ4本で、2つの図をならべます。

\begin{bkfigure}
\begin{tikzpicture}
\begin{axis}[bkplot, name=s1, width=56mm, height=46mm, xmin=0, xmax=1.15, ymin=0, ymax=1.15, xtick={0.25,0.5,0.75,1}, xticklabels={$\frac14$,$\frac24$,$\frac34$,1}, ytick={0.5,1}, xlabel={$x$}, ylabel={$y$}, title={\sffamily\small 0から1まで}]
  \foreach \k in {1,2,4} {\edef\tmp{\noexpand\fill[bkOrange!25, draw=bkOrange!70, line width=0.4pt] (axis cs:{(\k-1)/4},0) rectangle (axis cs:{\k/4},{(\k/4)^2});}\tmp}
  \fill[bkRed!40, draw=bkRed!80!black, line width=0.6pt] (axis cs:0.5,0) rectangle (axis cs:0.75,0.5625);
  \addplot[bkcurve, domain=0:1.05] {x^2};
\end{axis}
\begin{axis}[bkplot, at={(s1.east)}, anchor=west, xshift=16mm, width=56mm, height=46mm, xmin=0, xmax=2.3, ymin=0, ymax=4.6, xtick={0.5,1,1.5,2}, ytick={2,4}, xlabel={$x$}, ylabel={$y$}, title={\sffamily\small 0から2まで}]
  \foreach \k in {1,2,4} {\edef\tmp{\noexpand\fill[bkOrange!25, draw=bkOrange!70, line width=0.4pt] (axis cs:{(\k-1)/2},0) rectangle (axis cs:{\k/2},{(\k/2)^2});}\tmp}
  \fill[bkRed!40, draw=bkRed!80!black, line width=0.6pt] (axis cs:1,0) rectangle (axis cs:1.5,2.25);
  \addplot[bkcurve, domain=0:2.1] {x^2};
\end{axis}
\end{tikzpicture}
\bkfigcaption{2つの図は、形がまったく同じ。目もりの数字だけが、横は2倍、たては4倍になっている。}
\end{bkfigure}

赤い3本めの長方形をくらべます。

\begin{bkcenter}
\begin{bktable}{colspec={Q[l,m,wd=28mm] Q[c,m,wd=34mm] Q[c,m,wd=40mm]}, rowsep=2pt, row{1}={font=\sffamily\bfseries\small}}
赤い長方形 & 0から1まで & 0から2まで \\
はば & $1\div4=\dfrac{1}{4}$ & $2\div4=\dfrac{2}{4}$\quad（2倍） \\
右はし & $\dfrac{3}{4}$ & $\dfrac{2}{4}\times3=\dfrac{6}{4}$\quad（2倍） \\
高さ & $\left(\dfrac{3}{4}\right)^2=\dfrac{9}{16}$ & $\left(\dfrac{6}{4}\right)^2=\dfrac{36}{16}$\quad（4倍） \\
面積 & $\dfrac{9}{16}\times\dfrac{1}{4}=\dfrac{9}{64}$ & $\dfrac{36}{16}\times\dfrac{2}{4}=\dfrac{72}{64}$\quad（8倍） \\
\end{bktable}
\end{bkcenter}

はばが2倍になると、右はしの場所も2倍になります。高さは右はしの2乗なので、$2\times2=4$ 倍。面積は、はば2倍×高さ4倍で8倍です。ほかの長方形も同じなので、和も8倍、近づく先も8倍です。
\[ \int_0^2x^2\,\dd x=8\times\frac{1}{3}=\frac{8}{3} \]
8は $2^3$ です。0から $b$ までなら、はばが $b$ 倍、高さが $b^2$ 倍なので、面積は $b\times b^2=b^3$ 倍になります。

\begin{bkcenter}
\begin{bktable}{colspec={Q[l,m,wd=30mm] Q[c,m,wd=36mm] Q[c,m,wd=52mm]}, rowsep=2pt}
 & 0から1まで & 0から $b$ まで \\
はば & $\dfrac{1}{n}$ & $\dfrac{b}{n}$\quad（$b$ 倍） \\
$k$ 本めの右はし & $\dfrac{k}{n}$ & $\dfrac{b}{n}\times k=\dfrac{kb}{n}$\quad（$b$ 倍） \\
$k$ 本めの高さ & $\left(\dfrac{k}{n}\right)^2$ & $\left(\dfrac{kb}{n}\right)^2=b^2\left(\dfrac{k}{n}\right)^2$\quad（$b^2$ 倍） \\
長方形の面積 & $\left(\dfrac{k}{n}\right)^2\times\dfrac{1}{n}$ & $b\times b^2=b^3$ 倍 \\
\end{bktable}
\bktabcaption{はばも高さものびるので、どの長方形の面積も $b^3$ 倍になる。}
\end{bkcenter}
\[ \sum_{k=1}^{n}\left(\frac{kb}{n}\right)^2\times\frac{b}{n}=b^3\times\sum_{k=1}^{n}\left(\frac{k}{n}\right)^2\times\frac{1}{n}\ \longrightarrow\ b^3\times\frac{1}{3} \]
\noindent{\sffamily\bfseries\color{bkTheme}計算の結果}\quad $\displaystyle\int_0^b x^2\,\dd x=\frac{b^3}{3}$\par
\noindent $b=1$ なら $\dfrac{1}{3}$ で、④と合っています。$b=2$ なら $\dfrac{8}{3}$ で、上の計算と合っています。覚えなくても、「0から1までの $\dfrac{1}{3}$ を $b^3$ 倍」と考えれば、いつでも作れます。

\noindent{\sffamily\bfseries\color{bkTheme}0からでないとき}\par
0からでないときは、大きい面積から小さい面積をひきます。

\begin{bkfigure}
\begin{tikzpicture}
\pgfplotsset{mini/.style={bkplot, width=40mm, height=36mm, xmin=0, xmax=2.4, ymin=0, ymax=4.6, xtick={1,2}, xticklabels={$a$,$b$}, ytick=\empty}}
\begin{axis}[mini, name=m1]
  \addplot[fill=bkBlue!30, draw=none, domain=0:2] {x^2} \closedcycle;
  \addplot[bkcurve, line width=1.2pt, domain=0:2.1] {x^2};
\end{axis}
\node[font=\Large, anchor=west] (mi) at ([xshift=2mm]m1.east) {$-$};
\begin{axis}[mini, name=m2, at={(mi.east)}, anchor=west, xshift=2mm]
  \addplot[fill=bkGreen!35, draw=none, domain=0:1] {x^2} \closedcycle;
  \addplot[bkcurve, line width=1.2pt, domain=0:2.1] {x^2};
\end{axis}
\node[font=\Large, anchor=west] (eq) at ([xshift=2mm]m2.east) {$=$};
\begin{axis}[mini, name=m3, at={(eq.east)}, anchor=west, xshift=2mm]
  \addplot[fill=bkOrange!40, draw=none, domain=1:2] {x^2} \closedcycle;
  \addplot[bkcurve, line width=1.2pt, domain=0:2.1] {x^2};
\end{axis}
\end{tikzpicture}
\bkfigcaption{「0から $b$ まで」から「0から $a$ まで」をひくと、「$a$ から $b$ まで」が残る。}
\end{bkfigure}

\begin{align*}
\int_a^b x^2\,\dd x&=\int_0^b x^2\,\dd x-\int_0^a x^2\,\dd x\\
&=\frac{b^3}{3}-\frac{a^3}{3}
\end{align*}

\noindent{\sffamily\bfseries\color{bkTheme}グラフがまっすぐなとき}\par
グラフがまっすぐなら、5-9のように図形で求められます。$y=x$ の下は三角形、$y=1$ の下は長方形です。
\begin{align*}
\int_0^b x\,\dd x&=\frac{1}{2}\times b\times b=\frac{b^2}{2}&&\text{（三角形）}\\
\int_0^b 1\,\dd x&=1\times b=b&&\text{（長方形）}
\end{align*}
$y=x$ のときも、長方形の和で同じ答えになることを、あとの例題5で確かめます。

\begin{bkfigure}
\begin{tikzpicture}
\pgfplotsset{mini2/.style={bkplot, width=50mm, height=38mm, xmin=0, xmax=2.5, ymin=0, ymax=2.5, xtick={2}, xticklabels={$b$}, ytick=\empty, xlabel={$x$}, title style={xshift=3mm}}}
\begin{axis}[mini2, name=t1, title={\sffamily\small $y=x$ の下：三角形}]
  \addplot[fill=bkBlue!25, draw=none, domain=0:2] {x} \closedcycle;
  \addplot[bkcurve, line width=1.2pt, domain=0:2.3] {x};
  \node[font=\sffamily\footnotesize, anchor=west] at (axis cs:2.05,1.0) {高さ $b$};
\end{axis}
\begin{axis}[mini2, name=t2, at={(t1.east)}, anchor=west, xshift=22mm, title={\sffamily\small $y=1$ の下：長方形}]
  \addplot[fill=bkGreen!30, draw=none, domain=0:2] {1} \closedcycle;
  \addplot[bkcurve, line width=1.2pt, domain=0:2.3] {1};
  \node[font=\sffamily\footnotesize, anchor=south] at (axis cs:1,1.05) {高さ1};
\end{axis}
\end{tikzpicture}
\bkfigcaption{グラフがまっすぐなら、三角形（底辺 $b$、高さ $b$）や長方形（よこ $b$、たて1）の面積になる。}
\end{bkfigure}


\begin{senmonsho}[大学の教科書ふうの書き方]
\begin{bkhonbun}
区間 $[a,\,b]$ を $a=x_0<x_1<\cdots<x_n=b$ と分け、各区間 $[x_{k-1},\,x_k]$ から点 $\xi_k$ をとる。$f$ が連続ならば、和 $\displaystyle\sum_{k=1}^{n}f(\xi_k)\,(x_k-x_{k-1})$ は、分け方を細かくすると一定の値に収束する。これを $\displaystyle\int_a^b f(x)\,\dd x$ と書く。
\end{bkhonbun}
\bkyaku{$[a,\,b]$ は「$a$ から $b$ まで（両はしもふくむ）」。$\xi$（クシー）は、小さな区間の中から選んだ1点で、長方形の高さを決める場所。$x_k-x_{k-1}$ は $k$ 番めの長方形のはば $\Delta x$。この和を「リーマン和」とよぶ。}
\end{senmonsho}

\begin{taisetsu}
定積分\quad $\displaystyle\int_a^b f(x)\,\dd x$\\
＝ ① はば $\Delta x$ に分け、② 高さ $f(x)$×はばを1本ずつ計算し、③ $a$ から $b$ までたして、はばを0に近づけた先\\
＝ グラフと横軸のあいだの面積（横軸より下の部分はマイナス）\\
$\displaystyle\int_0^b x^2\,\dd x=\frac{b^3}{3}$\qquad $\displaystyle\int_0^b x\,\dd x=\frac{b^2}{2}$
\end{taisetsu}

\begin{kotoba}
\small
\bkword{$\Sigma$（シグマ）}{番号 $k$ を変えながら、全部たす記号。}
\bkword{定積分}{細かく分けてたした和の、はばを0に近づけた先。グラフと横軸のあいだの面積（下はマイナス）。}
\bkword{$\int$（インテグラル）}{定積分の記号。「たし合わせる」の S をのばした形。}
\bkword{下端・上端}{$\int_a^b$ の $a$（はじめ）と $b$（おわり）。}
\end{kotoba}

%% ---------------- 例題 ----------------
\bksection{reidai}{例題でたしかめよう}

\begin{reidai}[$\Sigma$ を計算する]
$\displaystyle\sum_{k=1}^{4}(2k^2-1)$ を計算しましょう。\par
\kangaekata $k$ に1、2、3、4を順に入れて、できた4つの数をたします。$2k^2$ は「$k$ を2乗してから2倍」です。\par
\shiki $k=1$：$2\times1^2-1=1$\qquad $k=2$：$2\times2^2-1=7$\par
\hspace*{2.6\zw}$k=3$：$2\times3^2-1=17$\qquad $k=4$：$2\times4^2-1=31$\par
\hspace*{2.6\zw}$1+7+17+31=56$\par
\kotae 56
\end{reidai}

\begin{reidai}[長方形の和を $\Sigma$ で書く]
$y=x^2$ のグラフの、$x=0$ から1までの下の面積を、右はしの長方形5本で近似します。和を $\Sigma$ で書き、計算しましょう。\par
\kangaekata ①の段1〜段4と同じです。はばは $1\div5=\dfrac{1}{5}$。$k$ 本めの右はしは、$\dfrac{1}{5}$ を $k$ こ進んだ $\dfrac{k}{5}$。高さは、そこでの $y=x^2$ の値 $\left(\dfrac{k}{5}\right)^2$ です。\par
\shiki $\displaystyle\sum_{k=1}^{5}\left(\frac{k}{5}\right)^2\times\frac{1}{5}=\sum_{k=1}^{5}\frac{k^2}{125}$\quad（$5^2\times5=125$）\par
\hspace*{2.6\zw}$\displaystyle=\frac{1+4+9+16+25}{125}=\frac{55}{125}=0.44$\par
\kotae 0.44（右はしなので、ほんとうの $\frac{1}{3}$ より大きい）
\end{reidai}

\noindent\begin{minipage}[t]{0.485\linewidth}\vspace{0pt}\raggedright
\begin{jibun}
$\displaystyle\sum_{k=1}^{5}k^2$ と $\displaystyle\sum_{k=1}^{3}(k+2)$ を計算しましょう。
\bkjans{55、12}
\end{jibun}
\end{minipage}\hfill
\begin{minipage}[t]{0.485\linewidth}\vspace{0pt}\raggedright
\begin{jibun}
同じ面積を、右はしの長方形3本で近似します。和を $\Sigma$ で書き、計算しましょう。
\bkjans{$\displaystyle\sum_{k=1}^{3}\left(\frac{k}{3}\right)^2\times\frac{1}{3}=\frac{1+4+9}{27}=\frac{14}{27}\approx0.52$}
\end{jibun}
\end{minipage}

\begin{reidai}[長方形の和で近似する]
$y=x$ のグラフの、$x=0$ から2までの下の面積を、はば0.5の右はしの長方形4本で近似しましょう。ほんとうの面積とくらべましょう。\par
\kangaekata 0から2までを、はば0.5で分けると $2\div0.5=4$ 本。右はしは0.5、1、1.5、2です。グラフは $y=x$ なので、高さは右はしの $x$ そのもの（2乗しない）です。\par
\shiki 高さは0.5、1、1.5、2。はばはどれも0.5なので、高さを先にたしてから0.5をかけます。\par
\hspace*{2.6\zw}$(0.5+1+1.5+2)\times0.5=5\times0.5=2.5$\par
\kotae 近似2.5（ほんとうは三角形で $\frac{1}{2}\times2\times2=2$。はみ出した分だけ大きい）
\end{reidai}

\begin{reidai}[$x^2$ の定積分]
$\displaystyle\int_0^3x^2\,\dd x$ を求めましょう。\par
\kangaekata $y=x^2$ のグラフの、0から3までの下の面積です。⑥で、0から $b$ までの面積は $\dfrac{b^3}{3}$ でした（0から1までの $\dfrac{1}{3}$ の $b^3$ 倍）。いまは上端が3なので、$b=3$ を入れます。\par
\shiki $\displaystyle\int_0^3x^2\,\dd x=\frac{3^3}{3}=\frac{27}{3}=9$\par
\bkstep[bkGreen]{\R{確}{たし}かめ}はば1の長方形なら、左はし $0+1+4=5$、右はし $1+4+9=14$。9は、そのあいだにあります。\par
\kotae 9
\end{reidai}

\begin{column}{1から100までの和}
240年ほど前のドイツで、小学生だったガウスは、先生に「1から100までを全部たしなさい」と言われました。ガウスはすぐに、5050と答えたそうです。1と100、2と99、3と98、……のように、はしから組にすると、どの組も和は101。組は $100\div2=50$ 組あるので、$101\times50=5050$。③の「階段を2つ組み合わせる」と同じ考え方で、式にすると $\dfrac{100\times101}{2}=5050$ です。ガウスは大人になって、数学だけでなく、電気や磁気の研究でも大きな仕事をしました。磁石のまわりの磁場の強さ（磁束密度）を表す単位「ガウス」は、彼の名前からとったものです（いまの国際単位ではテスラ（T）を使い、1\,T＝10000ガウスです）。
\end{column}

\noindent\begin{minipage}[t]{0.485\linewidth}\vspace{0pt}\raggedright
\begin{jibun}
$y=x^2$ のグラフの、$x=0$ から2までの下の面積を、はば0.5の右はしの長方形4本で近似しましょう。
\bkjans{$(0.25+1+2.25+4)\times0.5=3.75$（ほんとうは $\frac{8}{3}=2.67$）}
\end{jibun}
\end{minipage}\hfill
\begin{minipage}[t]{0.485\linewidth}\vspace{0pt}\raggedright
\begin{jibun}
$\displaystyle\int_0^2x^2\,\dd x$ と $\displaystyle\int_0^3x^2\,\dd x$ を求めましょう。
\bkjans{$\dfrac{2^3}{3}=\dfrac{8}{3}$、$\dfrac{3^3}{3}=9$}
\end{jibun}
\end{minipage}

\begin{reidai}[$y=x$ の下の面積を、長方形の和で求める]
$y=x$ のグラフの、$x=0$ から2までの下の面積を、右はしの長方形 $n$ 本の和で表し、$n$ をかぎりなく大きくして求めましょう。三角形の面積とくらべましょう。\par
\noindent\begin{minipage}[t]{0.66\linewidth}\raggedright
\kangaekata 段1〜段4で考えます。\par
\hspace*{2.6\zw}はば：長さ2を $n$ 等分して $\dfrac{2}{n}$\par
\hspace*{2.6\zw}右はし：$\dfrac{2}{n}$ を $k$ こ進んで $\dfrac{2k}{n}$\par
\hspace*{2.6\zw}高さ：$y=x$ なので、右はしと同じ $\dfrac{2k}{n}$\par
\hspace*{2.6\zw}面積：$\dfrac{2k}{n}\times\dfrac{2}{n}=\dfrac{4k}{n^2}$
\end{minipage}\hfill
\begin{minipage}[t]{0.32\linewidth}\centering\vspace{0pt}
\begin{tikzpicture}[x=15mm, y=15mm]
  \foreach \k in {1,...,4} {{\fill[bkOrange!30] ({(\k-1)/2},0) rectangle ({\k/2},{\k/2});
    \draw[bkOrange!80!black, line width=0.4pt] ({(\k-1)/2},0) rectangle ({\k/2},{\k/2});}}
  \draw[bkBlue!80!black, line width=1.1pt] (0,0) -- (2.1,2.1);
  \draw[bkInk, line width=0.6pt, -{Stealth[length=3pt]}] (0,0) -- (2.3,0) node[right, font=\footnotesize, inner sep=1pt] {$x$};
  \draw[bkInk, line width=0.6pt, -{Stealth[length=3pt]}] (0,0) -- (0,2.3) node[above, font=\footnotesize, inner sep=1pt] {$y$};
  \node[font=\footnotesize, anchor=north, inner sep=1pt] at (2,-0.05) {2};
  \node[font=\footnotesize, anchor=north, inner sep=1pt] at (1,-0.05) {1};
  \node[font=\footnotesize, text=bkBlue!80!black, anchor=south east, inner sep=1pt] at (1.3,1.55) {$y=x$};
\end{tikzpicture}\par
{\footnotesize\sffamily $n=4$ のとき}
\end{minipage}\par
\shiki $\displaystyle S_n=\sum_{k=1}^{n}\frac{4k}{n^2}=\frac{4}{n^2}\times(1+2+\cdots+n)$\par
\hspace*{2.6\zw}$\displaystyle\phantom{S_n}=\frac{4}{n^2}\times\frac{n(n+1)}{2}=\frac{2(n+1)}{n}=2\Bigl(1+\frac{1}{n}\Bigr)$\quad（③の和）\par
\hspace*{2.6\zw}$n=4$ なら $2\times1.25=2.5$（例題3と同じ）\par
\hspace*{2.6\zw}$n\to\infty$ で $\dfrac{1}{n}\to0$ なので、$2\times1=2$\par
\bkstep[bkGreen]{\R{確}{たし}かめ}三角形の面積は $\dfrac{1}{2}\times2\times2=2$。一致します。\par
\kotae 2
\end{reidai}

\begin{reidai}[曲がったグラフの距離]
速度が $v=3t^2$（m/s）の物体が、0秒から2秒までに進む距離を求めましょう。\par
\kangaekata 距離は $v$-$t$ グラフの面積なので、$\displaystyle\int_0^2 3t^2\,\dd t$ です。$v=3t^2$ のグラフは、$v=t^2$ をたてに3倍にのばした形です。長方形のはばはそのままで、高さがどれも3倍なので、面積の和も3倍です（下の図）。横軸の文字が $x$ から $t$ に変わっても形は同じなので、⑥の $\dfrac{b^3}{3}$ が使えます（$b=2$）。\par
\begin{center}
\begin{tikzpicture}[x=1mm, y=1mm]
  % t=1.25〜1.5 の細い長方形：灰（v=t^2 の高さ）と青（v=3t^2 の高さ）
  \fill[bkBlue!25] (25,0) rectangle (30,{3*1.5*1.5*2.6});
  \fill[bkGray!45] (25,0) rectangle (30,{1.5*1.5*2.6});
  \draw[bkGray!80!black, line width=0.4pt] (25,0) rectangle (30,{1.5*1.5*2.6});
  \draw[bkBlue!80!black, line width=0.4pt] (25,0) rectangle (30,{3*1.5*1.5*2.6});
  \draw[bkGray!80!black, line width=1.1pt] plot[domain=0:2, samples=30] ({20*\x},{2.6*\x*\x});
  \draw[bkBlue!80!black, line width=1.3pt] plot[domain=0:1.62, samples=30] ({20*\x},{7.8*\x*\x});
  \draw[bkInk, line width=0.6pt, -{Stealth[length=4pt]}] (0,0) -- (46,0) node[right, font=\sffamily\small, inner sep=1pt] {$t$};
  \draw[bkInk, line width=0.6pt, -{Stealth[length=4pt]}] (0,0) -- (0,23) node[above, font=\sffamily\small, inner sep=1pt] {$v$};
  \node[font=\sffamily\small, anchor=north] at (40,-0.5) {2};
  \node[font=\sffamily\small, anchor=north] at (30,-0.5) {1.5};
  \node[font=\sffamily\small, text=bkBlue!80!black, anchor=east] at (24,20) {$v=3t^2$};
  \node[font=\sffamily\small, text=bkGray!80!black, anchor=west] at (42,10.4) {$v=t^2$};
  % 高さの比べ（右に寸法線）
  \draw[bkInk, line width=0.5pt, <->] (32,0) -- (32,{1.5*1.5*2.6});
  \draw[bkBlue!80!black, line width=0.5pt, <->] (34.5,0) -- (34.5,{3*1.5*1.5*2.6});
  \node[font=\sffamily\footnotesize, anchor=west, inner sep=0pt] at (36,2.9) {1つ分};
  \node[font=\sffamily\bfseries\footnotesize, text=bkBlue!80!black, anchor=west, inner sep=0pt] at (36,17.5) {高さ3倍};
\end{tikzpicture}
\end{center}
\shiki $\displaystyle\begin{aligned}[t]\int_0^2 3t^2\,\dd t&=3\times\int_0^2 t^2\,\dd t\\&=3\times\frac{2^3}{3}=3\times\frac{8}{3}=8\end{aligned}$\par
\bkstep[bkGreen]{\R{確}{たし}かめ}はば1の長方形なら、左はし $0+3=3$、右はし $3+12=15$。8は、そのあいだにあります。単位は $(\text{m/s})\times\text{s}=\text{m}$ です。\par
\kotae \qty{8}{m}
\end{reidai}

\noindent\begin{minipage}[t]{0.485\linewidth}\vspace{0pt}\raggedright
\begin{jibun}
$y=x$ のグラフの、$x=0$ から3までの下の面積を、右はしの長方形 $n$ 本の和で表し、$n$ をかぎりなく大きくして求めましょう。三角形の面積とくらべましょう。
\bkjans{はば $\dfrac{3}{n}$、高さ $\dfrac{3k}{n}$。$\displaystyle\sum_{k=1}^{n}\frac{9k}{n^2}=\frac{9}{n^2}\times\frac{n(n+1)}{2}=\frac{9}{2}\Bigl(1+\frac{1}{n}\Bigr)\to\frac{9}{2}$。三角形も $\frac{1}{2}\times3\times3=\frac{9}{2}$}
\end{jibun}
\end{minipage}\hfill
\begin{minipage}[t]{0.485\linewidth}\vspace{0pt}\raggedright
\begin{jibun}
速度が $v=6t^2$（m/s）の物体が、0秒から3秒までに進む距離を求めましょう。
\bkjans{$6\times\dfrac{27}{3}=\qty{54}{m}$}
\end{jibun}
\end{minipage}

\begin{reidai}[とちゅうからの面積は、差で求める]
速度が $v=3t^2$（m/s）の物体が、1秒から3秒までに進む距離を求めましょう。\par
\kangaekata ⑥の「0からでないとき」と同じです。0秒から3秒までの面積から、0秒から1秒までの面積をひきます。高さが3倍なので、どちらも $t^2$ の面積の3倍です。\par
\shiki 0〜3秒：$\displaystyle3\times\frac{3^3}{3}=27$\qquad 0〜1秒：$\displaystyle3\times\frac{1^3}{3}=1$\par
\hspace*{2.6\zw}$27-1=26$\par
\kotae \qty{26}{m}
\end{reidai}

\begin{reidai}[台形で近似する]
$y=x^2$ の、$x=0$ から2までの下の面積を、はば0.5の台形4つで近似しましょう。\par
\kangaekata グラフ上の2点を結ぶ台形の面積は、（左の高さ＋右の高さ）÷2×はば（4-5）。左の高さだけで計算したのが左はしの長方形、右の高さだけで計算したのが右はしの長方形なので、台形の和は、2つの長方形の和の平均です（②の表の「2つの平均」）。\par
\shiki 左はし $(0+0.25+1+2.25)\times0.5=1.75$\par
\hspace*{2.6\zw}右はし $(0.25+1+2.25+4)\times0.5=3.75$\par
\hspace*{2.6\zw}平均 $\dfrac{1.75+3.75}{2}=2.75$\par
\kotae 2.75（ほんとうは $\frac{8}{3}=2.666\cdots$。長方形の3.75より、ずっと近い）
\end{reidai}

\noindent\begin{minipage}[t]{0.485\linewidth}\vspace{0pt}\raggedright
\begin{jibun}
速度が $v=t^2$（m/s）の物体が、2秒から4秒までに進む距離を求めましょう。
\bkjans{$\dfrac{64}{3}-\dfrac{8}{3}=\dfrac{56}{3}\approx\qty{18.7}{m}$}
\end{jibun}
\end{minipage}\hfill
\begin{minipage}[t]{0.485\linewidth}\vspace{0pt}\raggedright
\begin{jibun}
$y=x^2$ の、$x=0$ から1までの下の面積を、はば0.5の台形2つで近似しましょう。
\bkjans{左はし0.125、右はし0.625の平均で0.375（ほんとうは $\frac{1}{3}$）}
\end{jibun}
\end{minipage}

\begin{reidai}[うすい板に分けて体積を求める]
底面の半径 $r$、高さ $h$ の円すいの体積を求めましょう。\par
\kangaekata 頂点から $y$ の所で、底面と平行にうすく切ります（いちばん上の頂点で $y=0$、底面で $y=h$）。断面は円です。断面の半径は、頂点からの長さ $y$ に比例します（相似、4-3）。$y=h$ で半径 $r$ なので、$y$ の所では $r\times\dfrac{y}{h}=\dfrac{r}{h}y$。厚さ $\Delta y$ の円板（うすい円柱）を、上から下までたし合わせます。\par
\noindent\begin{minipage}[c]{0.62\linewidth}\raggedright
\shiki 円板の体積＝底面積×厚さ（4-5）\par
\hspace*{2.6\zw}$\pi\Bigl(\dfrac{r}{h}y\Bigr)^2\Delta y=\dfrac{\pi r^2}{h^2}\,y^2\,\Delta y$\par
\hspace*{2.6\zw}$\dfrac{\pi r^2}{h^2}$ は $y$ によらない数なので、外に出す\par
\hspace*{2.6\zw}$\displaystyle\int_0^h\frac{\pi r^2}{h^2}y^2\,\dd y=\frac{\pi r^2}{h^2}\times\frac{h^3}{3}=\frac{1}{3}\pi r^2h$\par
\kotae $\dfrac{1}{3}\pi r^2h$（4-5で学んだ公式と同じ）
\end{minipage}\hfill
\begin{minipage}[c]{0.36\linewidth}\centering
\begin{tikzpicture}[x=1mm, y=1mm]
  % 第6回：頂点 (0,32)、底面の中心 (0,2)、h=30、r=11。y は頂点から下向き。円板は y=18、厚さ Δy=2.4
  \def\yy{18}\def\dy{2.4}
  \pgfmathsetmacro\tp{32-\yy}\pgfmathsetmacro\bt{32-\yy-\dy}\pgfmathsetmacro\rr{11*\yy/30}
  \draw[bkInk, line width=0.5pt, dash pattern=on 1.5pt off 1pt] (11,2) arc[start angle=0, end angle=180, x radius=11, y radius=2.75];
  \fill[bkBlue!40] (-\rr,\tp) -- (-\rr,\bt) arc[start angle=180, end angle=360, x radius=\rr, y radius=\rr/4] -- (\rr,\tp) -- cycle;
  \fill[bkBlue!18] (0,\tp) ellipse [x radius=\rr, y radius=\rr/4];
  \draw[bkBlue!70!black, line width=0.6pt] (0,\tp) ellipse [x radius=\rr, y radius=\rr/4];
  \draw[bkBlue!70!black, line width=0.6pt] (-\rr,\tp) -- (-\rr,\bt) arc[start angle=180, end angle=360, x radius=\rr, y radius=\rr/4] -- (\rr,\tp);
  \draw[bkBlue!70!black, line width=0.6pt] (0,\tp) -- (\rr,\tp);
  \draw[bkInk, line width=0.7pt] (-11,2) arc[start angle=180, end angle=360, x radius=11, y radius=2.75];
  \draw[bkInk, line width=0.7pt] (-11,2) -- (0,32) -- (11,2);
  \draw[bkInk, line width=0.6pt] (0,2) -- (9.2,0.7);
  \fill[bkInk] (0,2) circle (0.5);
  \node[font=\sffamily\footnotesize, anchor=south, inner sep=1pt] at (4.6,1.4) {$r$};
  \draw[bkGray, line width=0.4pt] (-1,32) -- (-18.5,32) (-\rr-0.8,\tp) -- (-14.5,\tp) (-11.8,2) -- (-18.5,2);
  \draw[bkGray, line width=0.5pt, <->] (-13.5,32) -- (-13.5,\tp);
  \node[font=\sffamily\footnotesize, anchor=west] at (-13.1,{(32+\tp)/2}) {$y$};
  \draw[bkGray, line width=0.5pt, <->] (-17,32) -- (-17,2);
  \node[font=\sffamily\footnotesize, anchor=east] at (-17.3,17) {$h$};
  \draw[bkGray, line width=0.4pt] (\rr+0.8,\tp) -- (\rr+4.8,\tp) (\rr+0.8,\bt) -- (\rr+4.8,\bt);
  \draw[bkGray, line width=0.5pt, -{Stealth[length=3pt]}] ({\rr+3.8},{\tp+3}) -- ({\rr+3.8},\tp);
  \draw[bkGray, line width=0.5pt, -{Stealth[length=3pt]}] ({\rr+3.8},{\bt-3}) -- ({\rr+3.8},\bt);
  \node[font=\sffamily\footnotesize, anchor=west] at ({\rr+4.3},{(\tp+\bt)/2}) {$\Delta y$};
  \draw[bkBlue!70!black, line width=0.4pt] ({\rr/2},{\tp+0.2}) -- (6.6,23.4);
  \node[font=\sffamily\footnotesize, text=bkBlue!70!black, anchor=south west, align=left, inner sep=1pt] at (6.2,23.4) {半径\\$\dfrac{r}{h}y$};
\end{tikzpicture}
\end{minipage}
\end{reidai}

\begin{jibun}
底面の半径\qty{3}{cm}、高さ\qty{6}{cm}の円すいの体積を、\bkreidairef{9}と同じように積分で求めましょう。
\bkjans{$\displaystyle\int_0^6\frac{\pi\times9}{36}y^2\,\dd y=\frac{\pi}{4}\times72=18\pi\ \text{cm}^3$}
\end{jibun}

\begin{reidai}[左はしの長方形の和で、極限を求める]
$y=x^2$ のグラフの、$x=0$ から1までの下の面積を、左はしの長方形 $n$ 本の和で表し、$n$ をかぎりなく大きくして求めましょう。\par
\kangaekata 左はしの長方形では、$k$ 本めの高さは、左はし $\dfrac{k-1}{n}$ での値です。番号を1つずらして、0から $n-1$ までの $k$ で書くと、高さは $\left(\dfrac{k}{n}\right)^2$ です（まちがえやすいところ「$\Sigma$ の範囲」）。\par
\begin{center}
\begin{tikzpicture}[x=64mm, y=34mm]
  \foreach \k in {1,2,3} {
    \fill[bkBlue!25] ({\k/4},0) rectangle ({(\k+1)/4},{(\k/4)^2});
    \draw[bkBlue!70!black, line width=0.5pt] ({\k/4},0) rectangle ({(\k+1)/4},{(\k/4)^2});
}
  \draw[bkBlue!80!black, line width=1.2pt] plot[domain=0:1.04, samples=40] (\x,{\x*\x});
  \draw[bkInk, line width=0.6pt, -{Stealth[length=3pt]}] (0,0) -- (1.1,0) node[right, font=\footnotesize, inner sep=1pt] {$x$};
  \draw[bkInk, line width=0.6pt, -{Stealth[length=3pt]}] (0,0) -- (0,1.12) node[above, font=\footnotesize, inner sep=1pt] {$y$};
  \foreach \k in {0,1,2,3} {\node[font=\sffamily\footnotesize, anchor=north, inner sep=1pt] at ({(\k+0.5)/4},-0.04) {$k=\k$};
    \node[font=\footnotesize, anchor=north, inner sep=1pt] at ({(\k+0.5)/4},-0.2) {$\left(\frac{\k}{4}\right)^2$};}
  \node[font=\sffamily\footnotesize, anchor=east, inner sep=1pt] at (-0.02,-0.1) {番号};
  \node[font=\sffamily\footnotesize, anchor=east, inner sep=1pt] at (-0.02,-0.3) {高さ};
\end{tikzpicture}\par
{\footnotesize\sffamily $n=4$ の左はしの長方形。番号 $k$ を0から3にすると、高さは $\left(\frac{k}{4}\right)^2$}
\end{center}
\shiki $\displaystyle\sum_{k=0}^{n-1}\left(\frac{k}{n}\right)^2\times\frac{1}{n}=\sum_{k=1}^{n}\left(\frac{k}{n}\right)^2\times\frac{1}{n}-\left(\frac{n}{n}\right)^2\times\frac{1}{n}$\par
\hspace*{2.6\zw}（$k=0$ の項は0。右はしの和から、いちばん右の1本をひいた形）\par
\hspace*{2.6\zw}$\displaystyle=\frac{1}{6}\Bigl(1+\frac{1}{n}\Bigr)\Bigl(2+\frac{1}{n}\Bigr)-\frac{1}{n}$\quad（④の $S_n$）\par
\hspace*{2.6\zw}$n=4$ なら $0.46875-0.25=0.21875$（①と同じ）\par
\hspace*{2.6\zw}$n\to\infty$ で $\dfrac{1}{3}-0=\dfrac{1}{3}$\par
\kotae $\dfrac{1}{3}$（右はしの和と同じ値に近づく）
\end{reidai}

\begin{jibun}
$y=x$ のグラフの、$x=0$ から2までの下の面積を、左はしの長方形 $n$ 本の和で表し、$n$ をかぎりなく大きくして求めましょう。
\bkjans{右はしの和 $2\Bigl(1+\dfrac{1}{n}\Bigr)$（例題5）から、いちばん右の1本 $2\times\dfrac{2}{n}$ をひいて $2\Bigl(1-\dfrac{1}{n}\Bigr)\to2$}
\end{jibun}

\begin{column}{まん中の高さを使うと}
長方形の高さは、はばのまん中の値で決めてもかまいません（⑤の「せいかくに言うと」）。$y=x^2$ の0から1までを4本に分けると、まん中は $\frac{1}{8}$、$\frac{3}{8}$、$\frac{5}{8}$、$\frac{7}{8}$ なので、和は $\dfrac{1^2+3^2+5^2+7^2}{8^2}\times\dfrac{1}{4}=\dfrac{84}{256}=0.328125$。4本だけで、$\frac{1}{3}$ にかなり近くなります。はみ出す部分とすき間の部分が、ほぼ打ち消しあうからです。\par\vspace{2pt}

\noindent\begin{minipage}[c]{0.42\linewidth}\centering
\begin{tikzpicture}[x=1mm, y=1mm]
  % 1本の長方形を大きくした図。まん中の高さ h_m。左は すき間、右は はみ出し
  \fill[bkGreen!25] (0,0) rectangle (30,16);
  \fill[bkBlue!35] (0,16) -- (0,9.5) .. controls (8,11.6) and (12,13.6) .. (15,16) -- cycle;
  \fill[bkOrange!50] (15,16) .. controls (19,19) and (24,21.5) .. (30,24.5) -- (30,16) -- cycle;
  \draw[bkGreen!60!black, line width=0.5pt] (0,0) rectangle (30,16);
  \draw[bkBlue!80!black, line width=1.2pt] (-3,8.6) .. controls (-1,9) .. (0,9.5) .. controls (8,11.6) and (12,13.6) .. (15,16) .. controls (19,19) and (24,21.5) .. (30,24.5) -- (32,25.8);
  \draw[bkInk, line width=0.6pt] (-4,0) -- (34,0);
  \draw[bkGray, line width=0.4pt, dash pattern=on 1.2pt off 1pt] (15,0) -- (15,16);
  \node[font=\footnotesize, anchor=north, inner sep=1pt] at (15,-0.8) {まん中};
\end{tikzpicture}
\end{minipage}\hfill
\begin{minipage}[c]{0.55\linewidth}\raggedright
左の図は、1本の長方形を大きくしたものです。まん中より右では、グラフが長方形の上にはみ出しています（だいだい色）。左では、長方形の中にグラフより上の部分があり、余分にふくんでいます（青）。だいだい色と青の大きさがほぼ同じなので、打ち消しあって、和がほんとうの面積に近くなります。
\end{minipage}\par\vspace{2pt}
\noindent\begin{minipage}[c]{0.3\linewidth}\centering
\begin{tikzpicture}
\begin{axis}[bkplot, width=36mm, height=34mm, xmin=0, xmax=1.15, ymin=0, ymax=1.15, xtick={0.5,1}, ytick={0.5,1}, xlabel={$x$}, ylabel={$y$}]
  \foreach \k in {1,...,4} {\edef\tmp{\noexpand\fill[bkGreen!30, draw=bkGreen!70!black, line width=0.4pt] (axis cs:{(\k-1)/4},0) rectangle (axis cs:{\k/4},{((2*\k-1)/8)^2});}\tmp}
  \addplot[bkcurve, line width=1.2pt, domain=0:1.05] {x^2};
\end{axis}
\end{tikzpicture}
\end{minipage}\hfill
\begin{minipage}[c]{0.68\linewidth}\centering
\begin{bktable}{colspec={Q[c,m,wd=7mm] *{4}{Q[c,m,wd=15mm]}}, colsep=2pt, rowsep=1pt, cells={font=\footnotesize}, row{1}={font=\sffamily\bfseries\small}}
$n$ & 左はし & 右はし & 台形 & まん中 \\
4 & 0.21875 & 0.46875 & 0.34375 & 0.328125 \\
10 & 0.285 & 0.385 & 0.335 & 0.3325 \\
100 & 0.32835 & 0.33835 & 0.33335 & 0.333325 \\
\end{bktable}
\end{minipage}\par\vspace{2pt}
どの決め方でも、$n$ を大きくすると $\frac{1}{3}$ に近づきます。5-16では、まん中の高さを使って電気の量を計算します。
\end{column}

%% ---------------- まちがえやすいところ・物理工学 ----------------
\bksection{tsumazuki}{まちがえやすいところをチェック}

\begin{tsumazuki}{長方形の高さをまちがえる}
$y=x^2$ の0から1までを $n=4$ 本に分けたときの、$k$ 本めの高さ\par
\bkng{$k^2$（1、4、9、16）}
\bkng{$\dfrac{k}{4}$（2乗をわすれる）}
\bkok{$\left(\dfrac{k}{4}\right)^2$（$\frac{1}{16}$、$\frac{4}{16}$、$\frac{9}{16}$、$\frac{16}{16}$）}
$k$ は「何本め」という番号で、場所ではありません。右はしの場所は $\dfrac{k}{4}$ です。$k^2$ にすると、高さ16の長方形ができます。グラフは高さ1までしかないので、おかしいとわかります。高さは、グラフの式 $y=x^2$ に場所 $\dfrac{k}{4}$ を入れた値です。
\end{tsumazuki}

\begin{tsumazuki}{$b$ 倍にするのを、はばだけにする}
\bkng{$\displaystyle\int_0^2x^2\,\dd x=2\times\frac{1}{3}=\frac{2}{3}$}
\bkok{$\displaystyle\int_0^2x^2\,\dd x=2^3\times\frac{1}{3}=\frac{8}{3}$}
0から2までにすると、はばが2倍になるだけでなく、右はしが2倍になるので、高さも $2^2=4$ 倍になります。
\end{tsumazuki}

\begin{tsumazuki}{$\Sigma$ で、$k$ の範囲をまちがえる}
$\displaystyle\sum_{k=1}^{n}$ は $k=1$ から $n$ までの $n$ 個、$\displaystyle\sum_{k=0}^{n-1}$ は $k=0$ から $n-1$ までの $n$ 個です。左はしの長方形の和は $\displaystyle\sum_{k=0}^{n-1}\left(\frac{k}{n}\right)^2\frac{1}{n}$ で、右はしとは範囲が1つずれています。
\end{tsumazuki}

\begin{tsumazuki}{$\dd x$ をわすれる}
$\displaystyle\int_0^1x^2$ だけでは、何で細かく分けるのかがわかりません。$\dd x$ は「$x$ の方向に細く分けたはば」を表す大事な部分です。$\displaystyle\int v\,\dd t$ なら、時間で分けています。
\end{tsumazuki}

\begin{tsumazuki}{下端と上端を入れかえても同じだと思う}
\bkng{$\displaystyle\int_1^0x^2\,\dd x=\frac{1}{3}$}
\bkok{$\displaystyle\int_1^0x^2\,\dd x=-\frac{1}{3}$}
1から0へ分けていくと、はば $\Delta x$ が「あと${}-{}$まえ」でマイナスになります（⑤）。だから、符号が反対になります。
\end{tsumazuki}

\begin{butsuri}
\begin{itemize}[leftmargin=1.2\zw, itemsep=2pt, topsep=0pt, label={\color{bkPurple}●}]
\item \textbf{コンピュータの積分}\par
コンピュータは、式で積分できないときも、細い長方形や台形に分けてたし合わせて計算します。車の燃費計は、短い時間ごとの「燃料の流れる速さ×時間」をたして、使った燃料を求めています。
\item \textbf{CTスキャン}\par
CTスキャンでは、体をうすく輪切りにした断面の画像が得られます。断面の面積×厚さをたし合わせると、臓器などの体積がわかります（\bkreidairef{9}の円すいと同じ考え方）。
\item \textbf{面積計}\par
地図の上の形のまわりをなぞるだけで面積がわかる「プラニメーター」という道具は、積分の考え方で動いています。
\end{itemize}
\end{butsuri}


%% ---------------- 章末問題 ----------------
\bkendproblems
\bkproblemgroup{基本}{3}

\bkproblem 次を計算しましょう。\quad（1）$\displaystyle\sum_{k=1}^{6}k$\quad（2）$\displaystyle\sum_{k=1}^{4}3k$\quad（3）$\displaystyle\sum_{k=1}^{3}k^3$
\bkpans{（1）\nobreak 21\quad（2）\nobreak 30\quad（3）\nobreak 36}

\bkproblem $y=x^2$ のグラフの、$x=0$ から1までの下の面積を、右はしの長方形6本で近似します。和を $\Sigma$ で書き、計算しましょう（小数第3位まで）。
\bkpans{$\displaystyle\sum_{k=1}^{6}\left(\frac{k}{6}\right)^2\times\frac{1}{6}=\frac{91}{216}\approx0.421$}

\bkproblem $y=x^2$ のグラフの、$x=0$ から1までの下の面積を、はば0.5の左はしの長方形2つと、右はしの長方形2つで近似しましょう。
\bkpans{左はし：$(0+0.25)\times0.5=0.125$、右はし：$(0.25+1)\times0.5=0.625$}

\bkproblem 次の定積分を求めましょう。\quad（1）$\displaystyle\int_0^4x^2\,\dd x$\quad（2）$\displaystyle\int_0^5 x\,\dd x$
\bkpans{（1）\nobreak $\dfrac{64}{3}$\quad（2）\nobreak $\dfrac{25}{2}$}

\bkproblem 速度が $v=t^2$（m/s）の物体が、0秒から6秒までに進む距離を求めましょう。
\bkpans{$\dfrac{216}{3}=\qty{72}{m}$}

\bkproblem $\displaystyle\int_0^3x^2\,\dd x-\int_0^2x^2\,\dd x$ を計算しましょう。これは、どこの面積ですか。
\bkpans{$9-\dfrac{8}{3}=\dfrac{19}{3}$。$y=x^2$ の下の、$x=2$ から3までの面積}

\bkproblem 速度が $v=3t^2$（m/s）の物体が、1秒から2秒までに進む距離を求めましょう。
\bkpans{$8-1=\qty{7}{m}$}

\bkproblemgroup{チャレンジ}{3}

\bkproblem $y=x$ のグラフの下の面積を、長方形 $n$ 本の和で表し、$n$ をかぎりなく大きくして求めましょう。\quad（1）\nobreak $x=0$ から4まで、右はしの長方形\quad（2）\nobreak $x=0$ から3まで、左はしの長方形
\bkpans{（1）\nobreak $\displaystyle\sum_{k=1}^{n}\frac{4k}{n}\times\frac{4}{n}=\frac{16}{n^2}\times\frac{n(n+1)}{2}=8\Bigl(1+\frac{1}{n}\Bigr)\to8$\quad（2）\nobreak 右はしの和 $\dfrac{9}{2}\Bigl(1+\dfrac{1}{n}\Bigr)$ から、いちばん右の1本 $3\times\dfrac{3}{n}$ をひいて $\dfrac{9}{2}\Bigl(1-\dfrac{1}{n}\Bigr)\to\dfrac{9}{2}$}

\bkproblem 高さ $h$、底面が1辺 $a$ の正方形の四角すいの体積が $\dfrac{1}{3}a^2h$ になることを、\bkreidairef{9}と同じように積分で示しましょう。頂点から $y$ の所の断面は、1辺 $\dfrac{a}{h}y$ の正方形です（4-3）。
\bkpans{体積 $=\displaystyle\int_0^h\frac{a^2}{h^2}y^2\,\dd y=\frac{a^2}{h^2}\times\frac{h^3}{3}=\frac{1}{3}a^2h$}

\bkproblem $y=x^2$ の、$x=0$ から3までの下の面積を、はば1の台形3つで近似しましょう。ほんとうの値9とのずれは何\,\%ですか。
\bkpans{左はし $0+1+4=5$、右はし $1+4+9=14$、平均9.5。ずれ $\dfrac{0.5}{9}\approx5.6\,\%$}

\begin{furikaeri}
\begin{checklist}
\item 細い長方形の和で、面積を近似できる。\hfill{\footnotesize\sffamily\color{bkTheme}→\pageref{bk:5-10:hasami}ページ}
\item $\Sigma$ の記号を読み書きできる。\hfill{\footnotesize\sffamily\color{bkTheme}→\pageref{bk:5-10:sigma}ページ}
\item 定積分の意味を説明し、$x^2$ の定積分を求められる。\hfill{\footnotesize\sffamily\color{bkTheme}→\pageref{bk:5-10:teisekibun}ページ}
\end{checklist}
\end{furikaeri}

\bkfukidashi{delta}{0から $x$ までの面積は $\dfrac{x^3}{3}$。これを微分すると $x^2$ にもどるね。偶然かな？}
